Nach
[[Mannigfaltigkeit/Vektorbündel/R/Zusammenhang/Global integrabel/Globale Trivialisierung/Aufgabe|Aufgabe]]
liegt ein Isomorphismus von Vektorbündel vor. Nach
[[Mannigfaltigkeit/Vektorbündel/R/Zusammenhang/Horizontaler Schnitt/Vertikale Ableitung/Fakt|Fakt]]
gilt für einen horzontalen Schnitt $s$ die Beziehung

\mavergleichskette
{\vergleichskette
{ \nabla s
}
{ = }{ 0
}
{  }{
}
{    }{
}
{   }{
}
}
{}{}{.}
Auf einer offenen Menge

\mavergleichskette
{\vergleichskette
{ U
}
{ \cong }{ V
}
{ \subseteq }{  \R^n
}
{    }{
}
{   }{
}
}
{}{}{}
gilt daher insbesondere

\mavergleichskettedisp
{\vergleichskette
{ \nabla_{\partial_i} (s)
}
{ =} { 0
}
{ } {
}
{ } {
}
{ } {
}
}
{}{}{}
für alle $i$  und das bedeutet, dass sämtliche Christoffelsymbole der
\definitionsverweis {Basisschnitte}{}{}
$s_j$  bezüglich der Standardvektorfelder $\partial_i$  verschwinden.

 [[Kategorie:Latexseite]]
